\documentclass{article}
\usepackage[utf8]{inputenc}
\usepackage[spanish]{babel}
\usepackage{hyperref}
\usepackage{listings}
\usepackage{xcolor}
\usepackage{geometry}
\geometry{a4paper, margin=1.2in}

\title{Resumen de Asesoría: TechSphere Voice Agent}
\author{Sebastián Salazar Pérez}
\date{\today}

\begin{document}
\maketitle

\section{Problema Inicial}
El agente de voz desarrollado para el Tech Sphere challenge presentaba una voz muy robótica. Se evaluaron diferentes motores TTS conversacionales (como ElevenLabs o Cartesia), y se decidió utilizar \textbf{Gemini Multimodal (Audio Nativo)} para unificar el servicio y reducir la latencia.

\section{Solución Propuesta: Gemini Audio Nativo}
Se implementó la modalidad de audio nativo de Gemini (Audio-Out) a través de Google AI Studio. 

\subsection{Implementación en FastAPI (HTTP REST)}
Para integrar el audio nativo usando peticiones tradicionales, se configura el parámetro \texttt{response\_modalities=["AUDIO"]} en el cliente de GenAI.

\begin{lstlisting}[language=Python, backgroundcolor=\color{lightgray!20}, basicstyle=\ttfamily\footnotesize, breaklines=true, frame=single]
from google import genai
from google.genai import types
from fastapi import Response

client = genai.Client(api_key="TU_API_KEY_FREE_TIER")

def generar_respuesta_voz(contexto_rag: str, duda_paciente: str):
    prompt = f"Contexto Médico: {contexto_rag}\nPaciente: {duda_paciente}"
    
    response = client.models.generate_content(
        model='gemini-3.5-flash',
        contents=prompt,
        config=types.GenerateContentConfig(
            response_modalities=["AUDIO"],
            speech_config=types.SpeechConfig(
                voice_config=types.VoiceConfig(
                    prebuilt_voice_config=types.PrebuiltVoiceConfig(
                        voice_name="Aoede"
                    )
                )
            )
        )
    )
    
    audio_bytes = None
    for part in response.candidates[0].content.parts:
        if part.inline_data:
            audio_bytes = part.inline_data.data
            
    return audio_bytes
\end{lstlisting}

\section{Análisis de Cuotas y Límites (Free Tier)}
Al analizar el consumo de la capa gratuita (Free Tier) al habilitar la modalidad de voz, se determinó:
\begin{itemize}
    \item \textbf{Gemini 3.5 Flash:} Límite estricto de 20 peticiones por día (RPD). Se agota rápidamente en pruebas prolongadas.
    \item \textbf{Gemini 3.5 Flash Lite:} Límite más holgado de 500 RPD, pero menos robusto para la multimodalidad de audio directo.
    \item \textbf{Gemini 2.5 Flash Native Audio Dialog:} Permite llamadas ilimitadas al día y 1M TPM, diseñado explícitamente para WebSockets y streaming en vivo.
\end{itemize}

\section{Estrategia para el Demo}
Dado que el tiempo para la presentación ante el jurado es corto:
\begin{enumerate}
    \item \textbf{Si se mantiene HTTP (FastAPI):} Usar \texttt{gemini-3.5-flash} (o \texttt{gemini-2.5-flash}) pidiendo audio a través de \texttt{response\_modalities}, pero asegurándose de tener a la mano múltiples API Keys nuevas en un archivo \texttt{.env} para alternarlas rápidamente si se agota la cuota de 20 peticiones diarias durante la demostración en vivo.
    \item \textbf{Si se tienen WebSockets:} Migrar a la Live API de \texttt{Gemini 2.5 Flash Native Audio Dialog} usando el SDK asíncrono para lograr respuestas de latencia ultrabaja e ignorar las restricciones diarias de RPD de la versión 3.5 estándar.
\end{enumerate}

\end{document}